\section{SMART goal 1: Setup of an improved data quality report}
\subsection*{Goal definition}
\textbf{1) Setup of an improved data quality report}
    
\textbf{WHY:} This will help the team to run the daily checks faster

By the end of the implementation phase, the existing data quality report is to be updated to use the new infrastructure. This includes new reporting data pipelines and the new structure using star shema. Also, Power BI is to be used.

Success is measured by user feedback which has to be retrieved by \gls{ASTV} team. The feedback is to be implemented and the report must be on the productive systems.



\subsection*{Core results}
\begin{itemize}
    
\item \textbf{Update of data source:} In the implementation chapter, the data source for the third party checks was moved from beeing stored on a shared drive to beeing on sharepoint with a synchronisation onto \gls{AWSS3}. By this the data is versioned and available for other usage.

\item \textbf{Power BI Setup:} At the beginning of the implementation chapter, the setup of Power BI was performed. It used the existing guides of Swiss Post, and allowed in turn to start the development of the new improved \gls{DQM} report.

\item \textbf{Model data setup:} The data setup done trough out the course of the implementation chapters first section setup a model which was extentable. And allowed the usage of filters to be applied on all the data.

\item \textbf{New visuals to support faster interpretation by stakeholders:} The new visuals developed, allowed three things. The fast visual check if the data is correct on the four comparisions performed. They allow to see the changes on a timeline using the linechart and finally the table with the data allows the detailed check per day and location to be confirmed.

\item \textbf{Scrum review and after review feedback:} The scrum review showed the report to the whole \gls{ASTV} team and was a valuable means to get feedback. Further feedback was retrieved in separate meetings which was in turn applied. Resulting into the final version of the report beeing brought into production environment.

\item \textbf{New view Athena, and applied DevOps:} During the implementation of the report in Power BI, a separate view was setup in Athena. This new view allowed to use the new data sources while the old view was not touched and the two were not entangled. In addition for the deployment of this new view, the DevOps deployment process at Swiss Post was described and applied.
\end{itemize}




\subsection*{Supporting results}
The following results support the later implementation steps:
\begin{itemize}
    \item In the solution landscape chapter, the definition of Power BI as new reporting environment was established. Furthermore the change from wide table to star shema was established.

    \item In the solution architecture, the metrics are reflected in the definitions of the relevant submitted and discharged events, as described in Sections \ref{SolutionSubmittedEvent} and \ref{SolutionDischargedEvent}.
\end{itemize}

\subsection*{What was not achieved?}
All goals are fulfilled.

\subsection*{Goal verdict}
This goal was completely fulfilled. It has the report has been implemented including applying the feedback of the stakeholders.

\subsection*{Learning and outlook for future work}
In the importance of running checks between different environments from the first day of development. This is due to the issues that occured with the different environment of Power BI. Furthermore it's important to get the feedback of the stakeholders as soon as possible and plan enough time for the implementation of their requests.